\subsection{AUTOMORPHISMS AND REGULAR ELEMENTS}

Lemma 1. Let \(u\) be an automorphism of \(G\) , and \(H\) the set of its fixed points. a) \(H\) is a closed subgroup of \(G\) .

\begin{enumerate}
\def\labelenumi{\alph{enumi})}
\setcounter{enumi}{1}
\tightlist
\item
  If \(\mathrm{H}_0\) is central in G, then G is commutative (so \(\mathrm{G} = \mathrm{T}\) ).
\end{enumerate}

Assertion \(a\) ) is clear. To prove \(b\) ), we can replace \(\mathrm{G}\) by \(\mathrm{D}(\mathrm{G})\) ( \(\S 1\) , Cor. 1 of Prop. 4), and hence can assume that \(\mathrm{G}\) is semi-simple. Then, if \(\mathrm{H}_0\) is central in \(\mathrm{G}\) , we have \(\mathrm{L}(\mathrm{H}) = \{0\}\) , so the endomorphism \(\mathrm{L}(u) - \mathrm{Id}\) of \(\mathfrak{g}\) is bijective. Let \(f\) be the endomorphism of the manifold \(\mathrm{G}\) defined by \(f(g) = u(g)^{-1}g\) for \(g \in \mathrm{G}\) ; it is étale, for if \(g \in \mathrm{G}\) and \(x \in \mathfrak{g}\) , we have \(\mathrm{T}(f)(xg) = u(g)^{-1}(x - \mathrm{L}(u)(x))g\) , so the tangent map of \(f\) at \(g\) is bijective. It follows that the image of \(f\) is open and compact, hence coincides with \(\mathrm{G}\) since \(\mathrm{G}\) is connected. Now let \(\mathrm{E}\) be a framing of G (§4, no. 10, Def. 3) and \(u(\mathrm{E})\) its image under \(u\) . By Prop. 19 of loc. cit., there exists an element \(h\) of G such that \((\operatorname{Int} h)(\mathrm{E}) = u(\mathrm{E})\) . Let \(g \in \mathrm{G}\) be such that \(h = f(g) = u(g)^{-1}g\) ; then

\[
u \circ \operatorname{Int} g = (\operatorname{Int} u (g)) \circ u = \operatorname{Int} g \circ (\operatorname{Int} h) ^ {- 1} \circ u,
\]

so the framing \((\operatorname{Int} g)(\operatorname{E})\) is stable under u. If \((\operatorname{Int} g)(\operatorname{E}) = (\operatorname{T}_{1}, \operatorname{B}, (U_{\alpha})_{\alpha \in \operatorname{B}})\) , then \(\sum U_{\alpha} \in \operatorname{L}(\operatorname{H})\) ; since \(\operatorname{L}(\operatorname{H}) = \{0\}\) , this implies that \(B = \varnothing\) , so \(G = T_{1}\) , and G is commutative.

Lemma 2. Let x be an element of T and S a subtorus of T. If the identity component of \(Z(x) \cap Z(S)\) reduces to T, there exists an element s of S such that xs is regular.

For all \(\alpha\in\mathrm{R}(\mathrm{G},\mathrm{T})\) , let \(S_{\alpha}\) be the submanifold of S consisting of the elements s of S such that \((xs)^{\alpha}=1\) . If there is no element s of S such that xs is regular, S is the union of the submanifolds \(S_{\alpha}\) , hence is equal to one of them. Hence there exists \(\alpha\) in \(\mathrm{R}(\mathrm{G},\mathrm{T})\) such that \((xs)^{\alpha}=1\) for all \(s\in S\) ; but this implies that \(x^{\alpha}=1\) and \(\alpha|S=1\) , so \(\mathrm{Z}(x)\cap\mathrm{Z}(\mathrm{S})\supset\mathrm{Z}(\mathrm{Ker}\alpha)\) , hence the lemma.

Lemma 3. Assume that G is simply-connected. Let C be a chamber of t, and u an automorphism of G such that T and C are stable under u. Then the set of points of T fixed by u is connected.

Since G is simply-connected, \(\varGamma(\mathrm{T})\) is generated by the nodal vectors \(K_{\alpha}\) ( \(\alpha \in \mathrm{R}(\mathrm{G}, \mathrm{T})\) ), and hence has a basis consisting of the family of the \(K_{\alpha}\) for which \(\alpha\) belongs to the basis \(\mathrm{B}(\mathrm{C})\) defined by C (Chap. VI, §1, no. 10). Thus, it suffices to prove that, if \(\varphi\) is an automorphism of the torus T leaving stable a basis of \(\varGamma(\mathrm{T})\) , the set of fixed points of \(\varphi\) is connected. Decomposing this basis into the disjoint union of the orbits of the group generated by \(\varphi\) , we are reduced to the case in which \(T = U^{n}\) and \(\varphi\) is the automorphism \((z_{1}, \ldots, z_{n}) \mapsto (z_{2}, \ldots, z_{n}, z_{1})\) ; in this case the fixed points of \(\varphi\) are the points \((z, z, \ldots, z)\) for \(z \in U\) , which form a connected subgroup of T.

PROPOSITION 3. Let \(u\) be an automorphism of \(G\) , and let \(x\) be a point of \(G\) fixed by \(u\) .

\begin{enumerate}
\def\labelenumi{\alph{enumi})}
\item
  There exists an element \(a\) of \(\mathfrak{g}\) , fixed by \(\mathrm{L}(u)\) and by \(\operatorname{Ad} x\) , such that \(x \exp a\) is regular.
\item
  There exists a regular element \(g\) of \(G\) fixed by \(u\) and commuting with \(x\) .
\end{enumerate}

Let H be the group of fixed points of u, S a maximal torus of \(\mathrm{Z}(x)\cap\mathrm{H}\) , and K the identity component of \(\mathrm{Z}(\mathrm{S})\cap\mathrm{Z}(x)\) . This is a connected closed subgroup of G; further, by Cor. 5 of Th. 2 of §2, no. 2, there exist maximal tori of G containing S and x, so K is of maximal rank and contains S and x. On the other hand, K is stable under u since S and x are; denote by V the set of fixed points of u in K. Then

\[
\mathrm{S} \subset \mathrm{V} _ {0} \subset \mathrm{K} \cap \mathrm{H} \subset \mathrm{Z} (\mathrm{S}) \cap \mathrm{Z} (x) \cap \mathrm{H},
\]

so \(V_{0}\) is contained in the centralizer of S in \((\mathrm{Z}(x)\cap\mathrm{H})_{0}\) ; but the latter reduces to S (loc. cit., Cor. 6), hence finally \(V_{0}=S\) . Lemma 1 now implies that K is commutative, hence is a maximal torus of G (since it is connected and of maximal rank). It contains S and x, and is equal to the identity component of \(\mathrm{Z}(\mathrm{S})\cap\mathrm{Z}(x)\) ; assertion a) now follows from Lemma 2. We deduce b) by taking \(g=x\exp a\) .

COROLLARY. Let \(\mathfrak{s}\) be a compact Lie algebra, and let \(\varphi\) be an automorphism of \(\mathfrak{s}\) . There exists a regular element of \(\mathfrak{s}\) fixed by \(\varphi\) .

Replacing \(\mathfrak{s}\) by \(\mathcal{D}\mathfrak{s}\) , we can assume that \(\mathfrak{s}\) is semi-simple. Let \(S\) be a compact simply-connected Lie group with Lie algebra \(\mathfrak{s}\) , and let \(u\) be the automorphism of \(S\) such that \(L(u) = \varphi\) . Prop. 3 implies the existence of an element \(a\) of \(\mathfrak{s}\) , fixed by \(\varphi\) , such that \(\exp a\) is regular in \(S\) ; in particular, \(a\) is regular in \(\mathfrak{s}\) (no. 2, remark 4).

THEOREM 1. Let \(u\) be an automorphism of a connected compact Lie group \(G\) .

\begin{enumerate}
\def\labelenumi{\alph{enumi})}
\item
  The identity component of the group of fixed points of \(u\) contains a regular element of \(G\) .
\item
  There exists a maximal torus K of G and a chamber of L(K) that are stable under u.
\item
  If G is simply-connected, the set of fixed points of u is connected.
\end{enumerate}

Assertion a) is the particular case x = e of Prop. 3. We assume now that G is simply-connected and prove b) and c). Let x be an element of G fixed by u, and let g be a regular element of G, fixed by u and commuting with x (Prop. 3). The centralizer K of g is a maximal torus of G (no. 2, Remark 2), stable under u, and containing x and g. By Cor. 3 of Prop. 2 of no. 2, there exists a unique alcove A of L(K) such that \(g \in \exp A\) and \(0 \in \overline{A}\) ; since g is fixed by u, L(u) leaves A, and hence also the chamber of L(K) containing A, stable. This proves b); further, the set of points of K fixed by u is connected (Lemma 3) and contains x and e, hence c) (General Topology, Chap. I, §11, no. 1, Prop. 2).

It remains to prove \(b)\) in the general case. Now, if \(\tilde{\mathrm{D}}(\mathrm{G})\) is the universal covering of \(\mathrm{D}(\mathrm{G})\) , and if \(f: \tilde{\mathrm{D}}(\mathrm{G}) \to \mathrm{G}\) is the canonical morphism, there exists an automorphism \(\tilde{u}\) of \(\tilde{\mathrm{D}}(\mathrm{G})\) such that \(f \circ \tilde{u} = u \circ f\) . If \(\tilde{\mathrm{K}}\) is a maximal torus of \(\tilde{\mathrm{D}}(\mathrm{G})\) and \(\tilde{\mathrm{C}}\) a chamber of \(\mathrm{L}(\tilde{\mathrm{K}})\) , stable under \(\tilde{u}\) (this exists by what has already been proved), there exists (§2, no. 3, Remark 2) a unique maximal torus \(\mathrm{K}\) of \(\mathrm{G}\) and a unique chamber \(\mathrm{C}\) of \(\mathrm{L}(\mathrm{K})\) such that \(\tilde{\mathrm{K}} = f^{-1}(\mathrm{K})\) and \(\tilde{\mathrm{C}} = \mathrm{L}(f)^{-1}(\mathrm{C})\) , and we see immediately that \(\mathrm{K}\) and \(\mathrm{C}\) are stable under \(u\) , hence assertion \(b)\) in the general case.

COROLLARY 1. Assume that the Z-module \(\pi_{1}(G)\) is torsion-free.

\begin{enumerate}
\def\labelenumi{\alph{enumi})}
\item
  The centralizer of every element of G is connected.
\item
  Any two commuting elements of G belong to the same maximal torus.
\end{enumerate}

By Cor. 3 of Prop. 11 of §4, no. 6, D(G) is simply-connected. We have G = C(G) \(_0\) .D(G); let \(x \in\) G; write \(x = uv\) , with \(u \in\) C(G) \(_0\) and \(v \in\) D(G). Then Z(x) = C(G) \(_0\) .Z \(_{\text{D(G)}}\) (v). By Th. 1 c), Z \(_{\text{D(G)}}\) (v) is connected, so Z(x) is connected, hence a). Thus, by Cor. 3 of Th. 2 of §2, no. 2, Z(x) is the union of the maximal tori of G containing \(x\) , hence b).

COROLLARY 2. Let \(\Gamma\) be a compact subgroup of \(\operatorname{Aut}(\mathbf{G})\) with the following property:

(*) There exist elements \(u_{1}, \ldots, u_{n}\) of \(\Gamma\) such that, for all \(i\) , the closure \(\Gamma_{i}\) of the subgroup of \(\Gamma\) generated by \(u_{1}, \ldots, u_{i}\) is a normal subgroup of \(\Gamma\) , and \(\Gamma_{n} = \Gamma\) .

Then, there exists a maximal torus of G stable under the operation of \(\Gamma\) .

We argue by induction on the dimension of G. Clearly, we can assume that \(u_{1} \neq Id\) ; then the subgroup H of fixed points of \(u_{1}\) is distinct from G, and is stable under the operation of \(\varGamma\) . Moreover, since \(\varGamma\) is compact, the image of \(\varGamma\) in \(\mathrm{Aut}(\mathrm{H}_{0})\) is a quotient of \(\varGamma\) , hence it also satisfies condition (*). By the induction hypothesis, there exists a maximal torus S of H stable under \(\varGamma\) . The centralizer K of S in G is connected ( \(\S2\) , no. 2, Cor. 5) and stable under \(\varGamma\) ; this is a maximal torus of G, since \(H_{0}\) contains a regular element of G (Th. 1 a)) which is conjugate to an element of S (loc. cit., Cor. 4).

COROLLARY 3. Let H be a Lie group and \(\Gamma\) a compact subgroup of H. Assume that \(\mathrm{H}_0\) is compact and that \(\Gamma\) satisfies condition (*) of Cor. 2. Then there exists a maximal torus T of \(\mathrm{H}_0\) such that \(\Gamma \subset \mathrm{N}_{\mathrm{H}}(\mathrm{T})\) .

COROLLARY 4. Every nilpotent subgroup of a compact Lie group is contained in the normalizer of a maximal torus.

Let H be a compact Lie group, N a nilpotent subgroup of H. Then the closure \(\varGamma\) of N is also a nilpotent group (Chap. III, §9, no. 1, Cor. 2 of Prop. 1), and it suffices, in view of Cor. 3, to prove that \(\varGamma\) satisfies condition (*). Now \(\varGamma_{0}\) is a connected compact nilpotent Lie group, hence is a torus (§1, no. 4, Cor. 1 of Prop. 4), and there exists an element \(u_{1}\) of \(\varGamma\) generating a dense subgroup of \(\varGamma_{0}\) (General Topology, Chap. VII, §1, no. 3, text preceding Prop. 8). The finite group \(\varGamma/\varGamma_{0}\) is nilpotent and there exist \(\tilde{u}_{2},\ldots,\tilde{u}_{n}\in\varGamma/\varGamma_{0}\) generating \(\varGamma/\varGamma_{0}\) and such that the subgroup of \(\varGamma/\varGamma_{0}\) generated by \((\tilde{u}_{2},\ldots,\tilde{u}_{r})\) is normal for \(r=2,\ldots,n\) (Algebra, Chap. I, §6, no. 5, Th. 1 and no. 7, Th. 4). Then, if \(u_{2},\ldots,u_{n}\) are representatives of \(\tilde{u}_{2},\ldots,\tilde{u}_{n}\) in \(\varGamma\) , the sequence \((u_{1},\ldots,u_{n})\) has the required properties.

Example. Take \(\mathbf{G} = \mathbf{U}(n,\mathbf{C})\) . We shall see later that the subgroup of diagonal matrices in \(\mathbf{G}\) is a maximal torus of \(\mathbf{G}\) and that its normalizer is the set of monomial matrices (Algebra, Chap. II, §10, no. 7, Example II) in \(\mathbf{G}\) .

We conclude that, if \(\Phi\) is a separating positive hermitian form on a finite dimensional complex vector space V and \(\Gamma\) is a nilpotent subgroup of \(\mathbf{U}(\Phi)\) , there exists a basis of V for which the matrices of the elements of \(\Gamma\) are monomial (``Blichtfeldt's theorem'').

